\section{Alternativas}
Además de \LaTeX{}, existen alternativas para producir documentos con contenido matemático:

\begin{itemize}
    \item \href{https://typst.app/}{Typst}: alternativa moderna a \LaTeX{}, con una sintaxis más simple y una compilación considerablemente más rápida.
    \item Microsoft Word: Permite insertar una ecuación con el atajo \keys{\Alt + =}. Dentro de ese ambiente es posible escribir directamente en \LaTeX{}, seleccionando la opción \emph{LaTeX} en el menú de conversión \emph{Ecuación} (\emph{Equation} $\to$ \emph{Conversions} $\to$ \emph{LaTeX}), en lugar del formato lineal por defecto (UnicodeMath). Este último tiene el perjuicio de que dificulta seguir las convenciones internacionales de formato de escritura matemática.
\end{itemize}

Herramientas como GitHub, Markdown, Obsidian, Jupyter o cualquier motor basado en MathJax permiten escribir matemática usando sintaxis de \LaTeX{} y puede que en ellas sea más sencillo escribir el contenido no matemático y generar el documento final.
